
\subsection{Poincar\'e 代数}
Poincar\'e 变换的无穷小形式为：
\begin{equation}
    \Lambda^\mu_{\,\nu}=\delta^\mu_{\,\nu}+\omega^\mu_{\,\nu},\quad\,a^\mu=\epsilon^\mu
\end{equation}

这里,\begin{itemize}
    \item 由于保度规条件
    \[\begin{aligned}
        \eta_{\rho\sigma}&= \eta_{\rho\sigma}=\eta_{\mu\nu}\Lambda^\mu_{\;\;\rho}\Lambda^\nu_{\;\;\sigma}\\
        &=\eta_{\mu\nu}(\delta^\mu_{\;\;\rho}+\omega^\mu_{\;\;\rho})(\delta^\nu_{\;\;\sigma}+\omega^\nu_{\;\;\sigma})\\
        &=\eta_{\mu\nu}\delta^\mu_{\;\;\rho}\delta^\nu_{\;\;\sigma}+\eta_{\mu\nu}(\delta^\mu_{\;\;\rho}\omega^\nu_{\;\;\sigma}+\delta^\nu_{\;\;\sigma}\omega^\mu_{\;\;\rho})\\
        &=\eta_{\mu\nu}+\omega_{\sigma\rho}+\omega_{\rho\sigma}
    \end{aligned}\]
    我们得到\(\omega_{\mu\nu}\) 是反对称的。
    \item 由于 \(\omega_{\mu\nu}\) 是反对称的，它包含六个独立参数。再加上四个平移参数，这意味着
    \text{Poincar\'e 群} 一共具有十个参数。
\end{itemize}
\noindent
相应地，考虑，Poincar\'e 变换对应希尔伯特空间中的幺正算符的无穷小形式是：
\begin{equation}
    {U}(1+\omega,\epsilon)={1}+\frac{i}{2}\omega_{\rho\sigma}\mathscr{J}^{\rho\sigma}-i\epsilon_\rho \mathscr{P}^\rho+\ldots
\end{equation}
\begin{itemize}
    \item 显然，这里的生成元是厄米算符。
    \[\mathscr{J}^{\mu\nu\dagger}= \mathscr{J}^{\mu\nu},\quad \mathscr{P}^{\mu\dagger}=\mathscr{P}^\mu\]
\end{itemize}
\hrule
\vspace{2ex}
为了得到\(\mathscr{J}\) 和 \(\mathscr{P}\)代数关系;我们首先推导\(\mathscr{J}\) 和 \(\mathscr{P}\)在Poincar\'e变换下，如何变换：

\[ 
\begin{aligned}
        {U}(\Lambda,\,a){U}(1+\omega,\epsilon){U}^\dagger(\Lambda,\,a)&={U}(\Lambda,\,a){U}(1+\omega,\epsilon){U}(\Lambda^{-1},-\Lambda^{-1}\,a)\\
        &={U}(\Lambda,\,a){U}\left\{(1+\omega)\Lambda^{-1},-(1+\omega)\Lambda^{-1}\,a+\epsilon\right\}\\
        &={U}\left\{\Lambda(1+\omega)\Lambda^{-1},-\Lambda(1+\omega)\Lambda^{-1}\,a+\Lambda\epsilon+\,a\right\}\\
        &={U}\left\{\Lambda(1+\omega)\Lambda^{-1},\Lambda\epsilon-\Lambda\omega\Lambda^{-1}\,a\right\}
\end{aligned}
\]上式展开：
\[    \begin{aligned}
        L.H.S&={U}(\Lambda,\,a){U}(1+\omega,\epsilon){U}^\dagger(\Lambda,\,a)\\
        &={U}(\Lambda,\,a){U}\left\{{1}+\frac{i}{2}\omega_{\rho\sigma}\mathscr{J}^{\rho\sigma} -i\epsilon_\rho \mathscr{P}^\rho\right\}{U}^\dagger(\Lambda,\,a)\\
        &={1} +\frac{i}{2}\omega_{\rho\sigma}(\Lambda,\,a){U}\mathscr{J}^{\rho\sigma}{U}^\dagger(\Lambda,\,a)-i\epsilon_\rho{U}(\Lambda,\,a)\mathscr{P}^\rho{U}^\dagger(\Lambda,\,a)\\
    \end{aligned}\]
\[    \begin{aligned}
        R.H.S&={U}\left\{\Lambda(1+\omega)\Lambda^{-1},\Lambda\epsilon-\Lambda\omega\Lambda^{-1}\,a\right\}\\
        &={1} +\frac{i}{2}\left\{ \Lambda\omega\Lambda^{-1}\right\}_{\mu\nu}\mathscr{J}^{\mu\nu}-i\left\{\Lambda\epsilon-\Lambda\omega\Lambda^{-1}\,a\right\}_\mu\mathscr{P}^\mu\\
        &={1} +\frac{i}{2}\left\{ \Lambda_\mu^{\;\;\rho}\omega_{\rho\sigma}(\Lambda^{-1})^\sigma_{\,\nu}\mathscr{J}^{\mu\nu}\right\}-i\left\{\Lambda^{\;\;\rho}_\mu\epsilon_\rho\mathscr{P}^\mu-\Lambda_\mu^{\;\;\rho}\omega_{\rho\sigma}(\Lambda^{-1})^\sigma_{\,\nu}\,a^\nu \mathscr{P}^\mu\right\}\\
        &={1} +\frac{i}{2} \Lambda_\mu^{\;\;\rho}\omega_{\rho\sigma}\Lambda^{\;\;\sigma}_\nu\mathscr{J}^{\mu\nu}
        -i\Lambda^{\;\;\rho}_\mu\epsilon_\rho\mathscr{P}^\mu
        +i\Lambda_\mu^{\;\;\rho}\omega_{\rho\sigma}\Lambda^{\;\;\sigma}_\nu\,a^\nu\mathscr{P}^\mu\\
        %11
        &={1} +\frac{i}{2} \Lambda_\mu^{\;\;\rho}\omega_{\rho\sigma}\Lambda^{\;\;\sigma}_\nu\mathscr{J}^{\mu\nu}
        +\frac{i}{2}\left\{\Lambda_\mu^{\;\;\rho}\omega_{\rho\sigma}\Lambda^{\;\;\sigma}_\nu\,a^\nu\mathscr{P}^\mu+  
        \Lambda_\mu^{\;\;\sigma}\omega_{\sigma\rho}\Lambda^{\;\;\rho}_\nu\,a^\nu\mathscr{P}^\mu\right\}
        -i\Lambda^{\;\;\rho}_\mu\epsilon_\rho\mathscr{P}^\mu\\
        %22
        &={1} +\frac{i}{2}\omega_{\rho\sigma} \Lambda_\mu^{\;\;\rho}\Lambda^{\;\;\sigma}_\nu\mathscr{J}^{\mu\nu}
        +\frac{i}{2}\omega_{\rho\sigma}\left\{\Lambda_\mu^{\;\;\rho}\Lambda^{\;\;\sigma}_\nu\,a^\nu\mathscr{P}^\mu-
        \Lambda_\mu^{\;\;\sigma}\Lambda^{\;\;\rho}_\nu\,a^\nu\mathscr{P}^\mu\right\}
        -i\Lambda^{\;\;\rho}_\mu\epsilon_\rho\mathscr{P}^\mu\\
        %33
        &={1} +\frac{i}{2}\omega_{\rho\sigma} \left\{
        \Lambda_\mu^{\;\;\rho}\Lambda^{\;\;\sigma}_\nu\mathscr{J}^{\mu\nu}+
        \Lambda_\mu^{\;\;\rho}\Lambda^{\;\;\sigma}_\nu\,a^\nu\mathscr{P}^\mu-
        \Lambda_\mu^{\;\;\rho}\Lambda^{\;\;\sigma}_\nu\,a^\mu\mathscr{P}^\nu\right\}
        -i\Lambda^{\;\;\rho}_\mu\epsilon_\rho\mathscr{P}^\mu\\
        %44
        &={1} +\frac{i}{2}\omega_{\rho\sigma} \Lambda_\mu^{\;\;\rho}\Lambda^{\;\;\sigma}_\nu\left\{
        \mathscr{J}^{\mu\nu}+
        \,a^\nu\mathscr{P}^\mu-
        \,a^\mu\mathscr{P}^\nu\right\}
        -i\Lambda^{\;\;\rho}_\mu\epsilon_\rho\mathscr{P}^\mu\\
    \end{aligned}\]
我们得到
\[\Rightarrow \frac{i}{2}\omega_{\rho\sigma}{U}(\Lambda,\,a)\mathscr{J}^{\rho\sigma}{U}^\dagger(\Lambda,\,a)
    -i\epsilon_\rho{U}(\Lambda,\,a)\mathscr{P}^\rho{U}^\dagger(\Lambda,\,a)=
    \frac{i}{2}\omega_{\rho\sigma} \Lambda_\mu^{\;\;\rho}\Lambda^{\;\;\sigma}_\nu\left\{
    \mathscr{J}^{\mu\nu}+
    \,a^\nu\mathscr{P}^\mu-
    \,a^\mu\mathscr{P}^\nu\right\}
    -i\Lambda^{\;\;\rho}_\mu\epsilon_\rho\mathscr{P}^\mu
    \]
对比两边，我们得到 \(\mathscr{J}\) 和 \(\mathscr{P}\) 在Poincar\'e下的变换关系为：

        \begin{subequations}\label{4.3}
        \begin{align}
        &{U}(\Lambda,\,a)\mathscr{J}^{\rho\sigma}{U}^\dagger(\Lambda,\,a)
        =
        \Lambda_\mu^{\;\;\rho}\Lambda^{\;\;\sigma}_\nu
        \left(
        \mathscr{J}^{\mu\nu}
        +\,a^\nu\mathscr{P}^\mu
        -\,a^\mu\mathscr{P}^\nu
        \right)
        \label{4.3a}
        \\
        &{U}(\Lambda,\,a)\mathscr{P}^\rho{U}^\dagger(\Lambda,\,a)
        =
        \Lambda^{\;\;\rho}_\mu\mathscr{P}^\mu
        \label{4.3b}
    \end{align}
\end{subequations}
\hrule
\vspace{2ex}
接下来，我们对(\ref{4.3a})和(\ref{4.3b})两式的变换算符 \({U}(\Lambda,\,a)\rightarrow {U}(1+\omega,\epsilon)\)进行展开,首先是:
\[  {U}(\Lambda,\,a)\mathscr{J}^{\rho\sigma}{U}^\dagger(\Lambda,\,a)=
        \Lambda_\mu^{\;\;\rho}\Lambda^{\;\;\sigma}_\nu
        \left(
        \mathscr{J}^{\mu\nu}
        +\,a^\nu\mathscr{P}^\mu
        -\,a^\mu\mathscr{P}^\nu
        \right)\]

    \[
        \begin{aligned}
                L.H.S&={U}(\Lambda,\,a)\mathscr{J}^{\rho\sigma}{U}^\dagger(\Lambda,\,a)\\&={U}(1+\omega,\epsilon)\mathscr{J}^{\rho\sigma}{U}^\dagger(1+\omega,\epsilon)
                %1
                \\&=\left\{{1}+\frac{i}{2}\omega_{\mu\nu}\mathscr{J}^{\mu\nu}-i\epsilon_\mu \mathscr{P}^\mu+\ldots\right\}
                \mathscr{J}^{\rho\sigma}
                \left\{{1}+\frac{i}{2}\omega_{\mu\nu}\mathscr{J}^{\mu\nu}-i\epsilon_\mu \mathscr{P}^\mu+\ldots\right\}^\dagger
                %2
                \\&=\mathscr{J}^{\rho\sigma}
                +\frac{i}{2}\omega_{\mu\nu}\left\{\mathscr{J}^{\mu\nu}\mathscr{J}^{\rho\sigma}-\mathscr{J}^{\rho\sigma}\mathscr{J}^{\mu\nu} \right\}
                -i\epsilon_\mu\left\{\mathscr{P}^\mu\mathscr{J}^{\rho\sigma}-\mathscr{P}^\mu\mathscr{J}^{\rho\sigma} \right\}+\ldots
                %3
                \\&=\mathscr{J}^{\rho\sigma}
                +\frac{i}{2}\omega_{\mu\nu}\left[\mathscr{J}^{\mu\nu},\mathscr{J}^{\rho\sigma} \right]
                -i\epsilon_\mu\left[\mathscr{P}^\mu,\mathscr{J}^{\rho\sigma} \right]
    \end{aligned}
    \] 
    \[
        \begin{aligned}
                    R.H.S&=\Lambda_\mu^{\;\;\rho}\Lambda^{\;\;\sigma}_\nu
                    \left(\mathscr{J}^{\mu\nu}
                    +\,a^\nu\mathscr{P}^\mu
                    -\,a^\mu\mathscr{P}^\nu\right)
                    %1
                    \\&=\left(\delta_\mu^{\;\;\rho}+\omega_\mu^{\;\;\rho}\right)\left(\delta_\nu^{\;\;\sigma}+\omega_\nu^{\;\;\sigma}\right)
                    \left(\mathscr{J}^{\mu\nu}+\epsilon^\nu\mathscr{P}^\mu-\epsilon^\mu\mathscr{P}^\nu\right)
            %2
                    \\&=\mathscr{J}^{\rho\sigma}+
                    \delta_\mu^{\;\;\rho}\delta_\nu^{\;\;\sigma}\left\{\epsilon^\nu\mathscr{P}^\mu-\epsilon^\mu\mathscr{P}^\nu\right\}
                    +\left\{\delta_\mu^{\;\;\rho}\omega_\nu^{\;\;\sigma}+\delta_\nu^{\;\;\sigma}\omega_\mu^{\;\;\rho}\right\}\mathscr{J}^{\mu\nu}
                    +\left\{\delta_\mu^{\;\;\rho}\omega_\nu^{\;\;\sigma}+\delta_\nu^{\;\;\sigma}\omega_\mu^{\;\;\rho}\right\}
                    \left\{\epsilon^\nu\mathscr{P}^\mu-\epsilon^\mu\mathscr{P}^\nu\right\}
            %3
                    \\&=\mathscr{J}^{\rho\sigma}+
                    \left\{\delta_\mu^{\;\;\rho}\omega_\nu^{\;\;\sigma}\mathscr{J}^{\mu\nu}+\delta_\nu^{\;\;\sigma}\omega_\mu^{\;\;\rho}\mathscr{J}^{\mu\nu}\right\}
                    +\left\{\epsilon^\sigma\mathscr{P}^\rho-\epsilon^\rho\mathscr{P}^\sigma\right\}
                    +\left\{\delta_\mu^{\;\;\rho}\omega_\nu^{\;\;\sigma} \left[\epsilon^\nu\mathscr{P}^\mu-\epsilon^\mu\mathscr{P}^\nu\right]
                    +\delta_\nu^{\;\;\sigma}\omega_\mu^{\;\;\rho}\left[\epsilon^\nu\mathscr{P}^\mu-\epsilon^\mu\mathscr{P}^\nu\right]\right\}
            %4
                    \\&=\mathscr{J}^{\rho\sigma}+
                    \left\{\omega_\nu^{\;\;\sigma}\mathscr{J}^{\mu\sigma}+\omega_\mu^{\;\;\rho}\mathscr{J}^{\rho\nu}\right\}
                    +\left\{\epsilon_\mu\eta^{\mu\sigma}\mathscr{P}^\rho-\epsilon_\mu\eta^{\mu\rho}\mathscr{P}^\sigma\right\}
                    +\left\{\omega_\nu^{\;\;\sigma}\left[\epsilon^\nu\mathscr{P}^\rho-\epsilon^\rho\mathscr{P}^\nu\right]
                    +\omega_\mu^{\;\;\rho}\left[\epsilon^\sigma\mathscr{P}^\mu-\epsilon^\mu\mathscr{P}^\sigma\right]\right\}\\
            %5
                    &=\mathscr{J}^{\rho\sigma}+
                    \left\{\omega_{\mu\nu}\eta^{\nu\rho}\mathscr{J}^{\mu\sigma}+\omega_{\nu\mu}\eta^{\mu\sigma}\mathscr{J}^{\rho\nu}\right\}
                    +\epsilon_\mu\left\{\eta^{\mu\sigma}\mathscr{P}^\rho-\eta^{\mu\rho}\mathscr{P}^\sigma\right\}
                    \\&+\left\{\omega_\nu^{\;\;\sigma}\left[\epsilon^\nu\mathscr{P}^\rho-\epsilon^\rho\mathscr{P}^\nu\right]
                    +\omega_\nu^{\;\;\rho} \left[\epsilon^\rho\mathscr{P}^\nu-\epsilon^\nu\mathscr{P}^\rho\right]\right\}\\
                    &=\mathscr{J}^{\rho\sigma}+\frac{1}{2}
                    \left\{\omega_{\mu\nu}\eta^{\nu\rho}\mathscr{J}^{\mu\sigma}+ \omega_{\nu\mu}\eta^{\mu\rho}\mathscr{J}^{\nu\sigma}
                    +\omega_{\nu\mu}\eta^{\mu\sigma}\mathscr{J}^{\rho\nu}+\omega_{\mu\nu}\eta^{\nu\sigma}\mathscr{J}^{\rho\mu}\right\}
                    +\epsilon_\mu\left\{\eta^{\mu\sigma}\mathscr{P}^\rho-\eta^{\mu\rho}\mathscr{P}^\sigma\right\}\\
            %5
                    &=\mathscr{J}^{\rho\sigma}+\frac{1}{2}
                    \omega_{\mu\nu}\left\{\eta^{\nu\rho}\mathscr{J}^{\mu\sigma}-\eta^{\mu\rho}\mathscr{J}^{\nu\sigma}
                    -\eta^{\mu\sigma}\mathscr{J}^{\rho\nu}+\eta^{\nu\sigma}\mathscr{J}^{\rho\mu}\right\}
                    +\epsilon_\mu\left\{\eta^{\mu\sigma}\mathscr{P}^\rho-\eta^{\mu\rho}\mathscr{P}^\sigma\right\}
        \end{aligned}
    \]
我们得到:
\[
    \Rightarrow \frac{i}{2}\omega_{\mu\nu}\left[\mathscr{J}^{\mu\nu},\mathscr{J}^{\rho\sigma} \right]
    -i\epsilon_\mu\left[\mathscr{P}^\mu,\mathscr{J}^{\rho\sigma} \right]=
    \frac{1}{2}\omega_{\mu\nu}\left\{\eta^{\nu\rho}\mathscr{J}^{\mu\sigma}
    -\eta^{\mu\rho}\mathscr{J}^{\nu\sigma}
    -\eta^{\sigma\mu}\mathscr{J}^{\rho\nu}
    +\eta^{\sigma\nu}\mathscr{J}^{\rho\mu}\right\}
    +\epsilon_\mu\left\{\eta^{\mu\sigma}\mathscr{P}^\rho
    -\eta^{\mu\rho}\mathscr{P}^\sigma\right\} 
\]
对比两边，得到代数关系：
\begin{subequations}\label{4.4}
    \begin{align}
    &i\left[\mathscr{J}^{\mu\nu},\mathscr{J}^{\rho\sigma}\right]
    =\eta^{\nu\rho}\mathscr{J}^{\mu\sigma}
    -\eta^{\mu\rho}\mathscr{J}^{\nu\sigma}
    -\eta^{\sigma\mu}\mathscr{J}^{\rho\nu}
    +\eta^{\sigma\nu}\mathscr{J}^{\rho\mu}
    \\
    &i\left[\mathscr{P}^\mu,\mathscr{J}^{\rho\sigma}\right]
    =\eta^{\mu\rho}\mathscr{P}^\sigma
    -\eta^{\mu\sigma}\mathscr{P}^\rho
    \end{align}
\end{subequations}
\noindent
接着是\[{U}(\Lambda,\,a)\mathscr{P}^\rho{U}^\dagger(\Lambda,\,a)
        =\Lambda^{\;\;\rho}_\mu\mathscr{P}^\mu\] \[
        \begin{aligned}
                    L.H.S&={U}(\Lambda,\,a)\mathscr{P}^\rho{U}^\dagger(\Lambda,\,a)\\
    %1
                    &=\left\{{1}+\frac{i}{2}\omega_{\mu\nu}\mathscr{J}^{\mu\nu}-i\epsilon_\mu \mathscr{P}^\mu+\ldots\right\}
                    \mathscr{P}^\rho
                    \left\{{1}+\frac{i}{2}\omega_{\mu\nu}\mathscr{J}^{\mu\nu}-i\epsilon_\mu \mathscr{P}^\mu+\ldots\right\}^\dagger\\
    %2
                    &=\mathscr{P}^\rho+\frac{i}{2}\omega_{\mu\nu}\left\{ \mathscr{J}^{\mu\nu}\mathscr{P}^\rho-\mathscr{P}^\rho\mathscr{J}^{\mu\nu}\right\}
                     -i\epsilon_\mu\left\{ \mathscr{P}^\mu\mathscr{P}^\rho-\mathscr{P}^\rho\mathscr{P}^\mu\right\}\\    
    %3
                    &=\mathscr{P}^\rho+\frac{i}{2}\omega_{\mu\nu}\left[\mathscr{J}^{\mu\nu},\mathscr{P}^\rho\right]
                     -i\epsilon_\mu\left[ \mathscr{P}^\mu,\mathscr{P}^\rho\right]
        \end{aligned}
    \]
    \[
        \begin{aligned}
            R.H.S&=\Lambda^{\;\;\rho}_\mu\mathscr{P}^\mu
            \\&=\left\{\delta_\mu^{\;\;\rho}+\omega_\mu^{\;\;\rho}\right\}\mathscr{P}^\mu
            =\delta_\mu^{\;\;\rho}\mathscr{P}^\mu+\omega_\mu^{\;\;\rho}\mathscr{P}^\mu
            \\&=\mathscr{P}^\rho+\omega_{\mu\nu}\eta^{\nu\rho}\mathscr{P}^\mu
            =\mathscr{P}^\rho+\frac{1}{2}\omega_{\mu\nu}\left\{\eta^{\nu\rho}\mathscr{P}^\mu-\eta^{\mu\rho}\mathscr{P}^\nu\right\}
        \end{aligned}
    \]
得到：
\[
    \Rightarrow \frac{i}{2}\omega_{\mu\nu}\left[\mathscr{J}^{\mu\nu},\mathscr{P}^\rho\right]
    -i\epsilon_\mu\left[ \mathscr{P}^\mu,\mathscr{P}^\rho\right]=
    \frac{1}{2}\omega_{\mu\nu}\eta^{\nu\rho}\mathscr{P}^\mu-\eta^{\mu\rho}\mathscr{P}^\nu
\]

对比两边，得到代数关系：
\begin{subequations}\label{4.5}
        \begin{align}
        &\left[\mathscr{J}^{\mu\nu},\mathscr{P}^\rho\right]=
        \eta^{\nu\rho}\mathscr{P}^\mu-\eta^{\mu\rho}\mathscr{P}^\nu\\
        &\left[ \mathscr{P}^\mu,\mathscr{P}^\rho\right]=0
    \end{align}
\end{subequations}

\begin{tcolorbox}[colback=white!5!,colframe=black!45!]
 将(\refeq{4.4})(\refeq{4.5})写在一起,这些正是Poincar\'e群中生成元的代数关系：
        \begin{subequations}\label{4.6}
            \begin{align}
        &i\left[\mathscr{J}^{\mu\nu},\mathscr{J}^{\rho\sigma} \right]=
        \eta^{\nu\rho}\mathscr{J}^{\mu\sigma}
        -\eta^{\mu\rho}\mathscr{J}^{\nu\sigma}
        -\eta^{\sigma\mu}\mathscr{J}^{\rho\nu}
        +\eta^{\sigma\nu}\mathscr{J}^{\rho\mu}\\
        &i\left[\mathscr{P}^\mu,\mathscr{J}^{\rho\sigma} \right]=
        \eta^{\mu\rho}\mathscr{P}^\sigma-\eta^{\mu\sigma}\mathscr{P}^\rho\\
        &\left[ \mathscr{P}^\mu,\mathscr{P}^\rho\right]=0
    \end{align}
\end{subequations}
\end{tcolorbox}
    \vspace{2ex}
        \begin{center}
            \text{\Large ***}
        \end{center}
    \vspace{2ex}
在$0+1$维量子场论中，能量算符\(H = \mathscr{P}^{0}\)对易的算符扮是守恒的，是动量三矢以及角动量三矢

\[\mathbf{P} = \{\mathscr{P}^{1}, \mathscr{P}^{2}, \mathscr{P}^{3}\}\,,\\
\mathbf{J} = \{\mathscr{J}^{23}, \mathscr{J}^{31}, \mathscr{J}^{12}\}\]
当然，还有能量\(\mathscr{P}^{0}\)本身。剩下的生成元形成所谓的“boost”三矢
\[\mathbf{K} = \{\mathscr{J}^{01}, \mathscr{J}^{02}, \mathscr{J}^{03}\}\]
这些是不守恒的，这是为什么我们不使用K的本征值标记物理态的原因。3维中，对易关系(\refeq{4.6})可以重写为

\begin{subequations}
    \begin{align}
        [J_{i}, J_{j}] &= i\epsilon_{ijk} J_{k} , \\
        [J_{i}, K_{j}] &= i\epsilon_{ijk} K_{k} , \\
        [K_{i}, K_{j}] &= -i\epsilon_{ijk} J_{k} , \\
        [J_{i}, P_{j}] &= i\epsilon_{ijk} P_{k} , \\
        [K_{i}, P_{j}] &= -i H \delta_{ij} , \\
        [J_{i}, H] &= [P_{i}, H] = [H, H] = 0 , \\
        [K_{i}, H] &= -i P_{i} ,
    \end{align}
\end{subequations}
